% Sections 8.1 to 8.5, translated from sv/chapters/08/theory.tex.

\section{What is a function?}
\label{c8:sec:funktion}
A \textbf{function} $f$ is a rule that gives exactly one value $f(x)$ for every permitted value of
$x$:
\[
  f : x \mapsto f(x)
\]
Functions can be represented as algebraic expressions, tables or graphs.

\begin{exempel}
$f(x) = 2x + 1$ gives $f(3) = 7$ and $f(-1) = -1$.
\end{exempel}

\section{Domain and range}
\label{c8:sec:mangder}
The \textbf{domain} $D_f$ is the set of all $x$-values for which $f(x)$ is defined.

The \textbf{range} $V_f$ is the set of all possible function values $f(x)$.

\begin{narrowtable}{@{}p{7em}p{8.5em}p{8.5em}@{}}
  \thead{Function} & \thead{Domain} & \thead{Range} \\
  \midrule
  $f(x) = 2x + 3$ & $\mathbb{R}$ & $\mathbb{R}$ \\
  $f(x) = x^2$ & $\mathbb{R}$ & $[0, \infty)$ \\
  $f(x) = \sqrt{x}$ & $[0, \infty)$ & $[0, \infty)$ \\
  $f(x) = \dfrac{1}{x}$ & $\mathbb{R} \setminus \{0\}$ & $\mathbb{R} \setminus \{0\}$ \\
\end{narrowtable}

\section{Common types of function}
\label{c8:sec:funktionstyper}

\subsection{Linear function}
\label{c8:sec:linjar}
\[
  f(x) = kx + m
\]
where $k$ is the slope and $m$ is the $y$-intercept.

\subsection{Power function}
\label{c8:sec:potens}
\[
  f(x) = x^n
\]

\subsection{Exponential function}
\label{c8:sec:exponential}
\[
  f(x) = C \cdot a^x, \qquad a > 0,\ a \neq 1
\]
\begin{itemize}
  \item $a > 1$: an increasing function.
  \item $0 < a < 1$: a decreasing function.
  \item $f(0) = C$, the initial value.
\end{itemize}

\subsection{Polynomial function}
\label{c8:sec:polynom}
\[
  f(x) = a_n x^n + a_{n-1} x^{n-1} + \cdots + a_0
\]

\section{Reading a graph}
\label{c8:sec:graf}
From a graph you can read:
\begin{itemize}
  \item the domain, that is, which $x$-values exist,
  \item the range, that is, which $y$-values are reached,
  \item the zeros, the $x$ where $f(x) = 0$,
  \item the maximum and minimum points.
\end{itemize}

\section{Worked examples}
\label{c8:sec:typexempel}

\begin{exempel}[Temperature sensor]
A sensor gives the temperature $f(x)$ (in $^{\circ}$C) as a function of the input voltage $x$ (in
V):
\[
  f(x) = 100x - 50, \qquad 0 \leq x \leq 5\,\text{V}
\]
\textbf{a)} Find the domain and the range. The domain is $D_f = [0, 5]\,\text{V}$, and
\[
  f(0) = -50\,^{\circ}\text{C}, \quad f(5) = 450\,^{\circ}\text{C}
  \quad \Rightarrow \quad V_f = [-50, 450]\,^{\circ}\text{C}.
\]
\textbf{b)} Calculate the input voltage when the temperature is $30\,^{\circ}\text{C}$.
\[
  100x - 50 = 30 \quad \Rightarrow \quad x = 0.8\,\text{V}
\]
\end{exempel}

\begin{exempel}[A linear function from two points]
A linear function passes through $(1, 3)$ and $(3, 7)$. The slope is
\[
  k = \frac{7 - 3}{3 - 1} = 2.
\]
Substituting $(1, 3)$ gives $m = 3 - 2 \cdot 1 = 1$, so
\[
  f(x) = 2x + 1.
\]
\end{exempel}

\begin{exempel}[Exponential function and growth]
A data set starts with $6 \times 10^6$ data points and grows by $3\,\%$ a year:
\[
  f(x) = 6 \times 10^6 \cdot 1.03^x
\]
\textbf{a)} Calculate $f(5)$.
\[
  f(5) = 6 \times 10^6 \cdot 1.03^5 \approx 6.96 \times 10^6
\]
\textbf{b)} When has the data set doubled?
\[
  1.03^x = 2 \quad \Rightarrow \quad x = \frac{\log 2}{\log 1.03} \approx 23.5\,\text{years}
\]
\end{exempel}
